%-------------------------
% Resume template — Jake's-resume style, content stripped.
% Placeholders use {field} / {{#each list}} matching the profile schema:
%   header: {name} {phone} {email} {linkedin} {github}
%   sections in THIS order: Education, Experience, Projects, Skills.
% Skills: group profile skills[] into Languages vs Frameworks/Tools when
% possible; otherwise one line. Keep the output to ONE page.
%-------------------------
\documentclass[letterpaper,11pt]{article}

\usepackage{latexsym}
\usepackage[empty]{fullpage}
\usepackage{titlesec}
\usepackage{marvosym}
\usepackage[usenames,dvipsnames]{color}
\usepackage{verbatim}
\usepackage{enumitem}
\usepackage[hidelinks]{hyperref}
\usepackage{fancyhdr}
\usepackage[english]{babel}
\usepackage{tabularx}
% glyphtounicode / \pdfgentounicode are pdfTeX-only and make the PDF's text
% extractable for ATS. XeTeX (tectonic) emits Unicode-mapped text natively and
% errors on them, so guard both — the template must build on either engine.
\usepackage{iftex}
\ifPDFTeX\input{glyphtounicode}\fi

\pagestyle{fancy}
\fancyhf{}
\fancyfoot{}
\renewcommand{\headrulewidth}{0pt}
\renewcommand{\footrulewidth}{0pt}

\addtolength{\oddsidemargin}{-0.5in}
\addtolength{\evensidemargin}{-0.5in}
\addtolength{\textwidth}{1in}
\addtolength{\topmargin}{-.5in}
\addtolength{\textheight}{1.0in}

\urlstyle{same}

\raggedbottom
\raggedright
\setlength{\tabcolsep}{0in}

\titleformat{\section}{
\vspace{-4pt}\scshape\raggedright\large
}{}{0em}{}[\color{black}\titlerule \vspace{-5pt}]

\ifPDFTeX\pdfgentounicode=1\fi

% Two arguments, per the campaign contract: an optional bold label and the
% bullet body. The assembler always emits \resumeItem{}{...}; the empty-label
% test keeps that case byte-identical to a plain bullet, with no stray indent.
\newcommand{\resumeItem}[2]{
\item\small{
{\ifx\relax#1\relax\else\textbf{#1}\hspace{0.5mm}\fi#2 \vspace{-1pt}}
}
}

\newcommand{\resumeSubheading}[4]{
\vspace{-2pt}\item
\begin{tabular*}{0.97\textwidth}[t]{l@{\extracolsep{\fill}}r}
\textbf{#1} & #2 \\
\textit{\small#3} & \textit{\small #4} \\
\end{tabular*}\vspace{-7pt}
}

\newcommand{\resumeSubSubheading}[2]{
\item
\begin{tabular*}{0.97\textwidth}{l@{\extracolsep{\fill}}r}
\textit{\small#1} & \textit{\small #2} \\
\end{tabular*}\vspace{-7pt}
}

\newcommand{\resumeProjectHeading}[2]{
\item
\begin{tabular*}{0.97\textwidth}{l@{\extracolsep{\fill}}r}
\small#1 & \textit{\small #2} \\
\end{tabular*}\vspace{-5pt}
}

\newcommand{\resumeSubItem}[2]{\resumeItem{#1}{#2}\vspace{-4pt}}

\renewcommand\labelitemii{$\vcenter{\hbox{\tiny$\bullet$}}$}

\newcommand{\resumeSubHeadingListStart}{\begin{itemize}[leftmargin=0.15in, label={}]}
\newcommand{\resumeSubHeadingListEnd}{\end{itemize}}
\newcommand{\resumeItemListStart}{\begin{itemize}}
\newcommand{\resumeItemListEnd}{\end{itemize}\vspace{-5pt}}

% The Skills section is a label-less list of "Category: a, b, c" rows, and the
% transition into Projects has to stay tight enough for the gap the machine
% gate measures between the last work line and the Projects heading.
\newcommand{\resumeHeadingSkillStart}{\begin{itemize}[leftmargin=0.15in, label={}]}
\newcommand{\resumeHeadingSkillEnd}{\end{itemize}\vspace{-5pt}}
\newcommand{\resumeSectionTransition}{\vspace{-9mm}}

\begin{document}

\begin{center}
\textbf{\Huge \scshape {name}} \\ \vspace{1pt}
\small
\href{sms:{phone}}{{phone}} $|$
\href{mailto:{email}}{{email}} $|$
\href{https://www.linkedin.com/in/{linkedin}}{\underline{linkedin.com/in/{linkedin}}} $|$
\href{https://github.com/{github}}{\underline{github.com/{github}}}
\end{center}

% Section order, spacing and scaffolding are contract, not decoration: the
% campaign assembler reads the \vspace before each \section, the indent between
% \resumeSubHeadingListStart and \resumeSubheading, and the Projects tail
% straight out of this file, then emits a resume shaped to match. Sections are
% Education / Skills / Working Experience / Projects, in that order; every
% bullet takes the two-argument form: an optional bold label first, then the
% bullet body. Each section shows one literal entry as the example, and the
% entries must not be wrapped in loops: the scaffolding is measured literally.
\section{\textbf{Education}}
\vspace{-0.4mm}
\resumeSubHeadingListStart
  \resumeSubheading
    {institution}
    {date}
    {degree}
    {location}
\resumeSubHeadingListEnd

\vspace{-4mm}
\section{\textbf{Skills}}
\resumeHeadingSkillStart
\resumeSubItem{Programming Languages:}
  {languages-from-skills}
\vspace{-1mm}
\resumeSubItem{Tools \& Technologies:}
  {frameworks-and-tools-from-skills}
\resumeHeadingSkillEnd

\vspace{-4mm}
\section{\textbf{Working Experience}}
\resumeSubHeadingListStart
  \resumeSubheading
    {company}
    {date}
    {title}
    {location}
    \resumeItemListStart
      \resumeItem{}{experience-bullet}
    \resumeItemListEnd
\resumeSubHeadingListEnd

\resumeSectionTransition
\section{\textbf{Projects}}
\vspace{-1mm}
\resumeSubHeadingListStart
  \resumeSubheading
    {project-name}
    {technologies}
    {}
    {date-range}
    \resumeItemListStart
      \resumeItem{}{project-bullet}
    \resumeItemListEnd
\resumeSubHeadingListEnd
\vspace{-4mm}
\resumeSubHeadingListStart
  \resumeSubheading
    {second-project-name}
    {technologies}
    {}
    {date-range}
    \resumeItemListStart
      \resumeItem{}{project-bullet}
    \resumeItemListEnd
\resumeSubHeadingListEnd
\vspace{-4mm}

\end{document}
